\chapter{Diagrammes FAST}

Le diagramme \textbf{FAST} se lit de gauche à droite avec la question \og \emph{comment} ? \fg{} ; dans l'autre sens, la même chaîne répond à \og \emph{pourquoi} ? \fg. Chaque sous-fonction est reliée, par une flèche pointillée, à la solution technique (en jaune) qui la réalise. Dans tout le document, $W$ désigne la matrice de connexions du réseau.

\noindent{\small\textbf{Légende.}~%
\tikz\fill[whybg, draw=whyborder, line width=0.5pt] (0,0) rectangle (0.26,0.13);~fonction \quad
\tikz\fill[gray!8, draw=cardline, line width=0.5pt] (0,0) rectangle (0.26,0.13);~sous-fonction \quad
\tikz\fill[bmMid, draw=pocborder, line width=0.5pt] (0,0) rectangle (0.26,0.13);~solution technique.}


\begin{figure}[H]
\centering
\scalebox{0.75}{%
\begin{tikzpicture}[fast]
% Colonnes : fonction (x=0), bus vertical (x=2.4), sous-fonctions (x=6.6), solutions (x=12.6)
\def\xbus{2.4}
\node[fct] (FP1) at (0,0) {FP1\\[2pt]Simuler des réseaux\\Hodgkin-Huxley\\couplés par synapses};

\node[sous] (S1) at (6.6, 3.6) {Décrire le réseau à simuler};
\node[sous] (S2) at (6.6, 1.2) {Faire communiquer les neurones entre eux};
\node[sous] (S3) at (6.6,-1.2) {Faire évoluer le potentiel de chaque neurone};
\node[sous] (S4) at (6.6,-3.6) {Enchaîner les pas de temps jusqu'à la durée};

\node[sol] (T1) at (12.6, 3.6) {Générateur de connectome :\\taille $N$, chaîne,\\tout-à-tout, aléatoire (FS1)\\(\textit{conf.json} pour description du réseau\\Python pour générer la matrice $W$)};
\node[sol] (T2) at (12.6, 1.2) {Chaque neurone intègre\\les messages de ses voisins, selon $W$};
\node[sol] (T3) at (12.6,-1.2) {Résolution Hodgkin-Huxley,\\équations du socle (FC1)};
\node[sol] (T4) at (12.6,-3.6) {Boucle de simulation,\\pas de temps du socle};

\node[font=\scriptsize\itshape, text=cardink] at (6.6,5.0) {Comment ? $\rightarrow$ \quad $\leftarrow$ Pourquoi ?};

% Arborescence orthogonale : tronc commun puis une branche par sous-fonction
\draw[tronc] (FP1.east) -- (\xbus,0);
\draw[tronc] (\xbus,3.6) -- (\xbus,-3.6);
\node[noeud] at (\xbus,0) {};
\foreach \i in {1,...,4}{
    \draw[lien] (\xbus,0 |- S\i.west) -- node[cmt]{Comment ?} (S\i.west);
    \draw[real] (S\i.east) -- (T\i.west);
}
\end{tikzpicture}}%
\caption{FAST de FP1 : le réseau est décrit par le générateur de connectome (FS1), puis simulé pas à pas -- communication entre les neurones, mise à jour Hodgkin-Huxley (FC1) -- jusqu'à la durée choisie.}
\label{fig:fast-fp1}
\end{figure}

\begin{figure}[H]
\centering
\scalebox{0.75}{%
\begin{tikzpicture}[fast]
\def\xbus{2.4}
\node[fct] (FP2) at (0,0) {FP2\\[2pt]Accélérer la\\simulation\\sur GPU};

\node[sous] (S1) at (6.6, 2.4) {Répartir les neurones sur les c\oe urs};
\node[sous] (S2) at (6.6, 0)   {Garder les données sur la carte};
\node[sous] (S3) at (6.6,-2.4) {Piloter la carte sans la ralentir};

\node[sol] (T1) at (12.6, 2.4) {Un fil de calcul par neurone};
\node[sol] (T2) at (12.6, 0)   {VRAM résidente,\\transferts début et fin (FC2)};
\node[sol] (T3) at (12.6,-2.4) {Orchestrateur CPU (Python) :\\lance les calculs à chaque pas};

\node[font=\scriptsize\itshape, text=cardink] at (6.6,3.8) {Comment ? $\rightarrow$ \quad $\leftarrow$ Pourquoi ?};

\draw[tronc] (FP2.east) -- (\xbus,0);
\draw[tronc] (\xbus,2.4) -- (\xbus,-2.4);
\node[noeud] at (\xbus,0) {};
\foreach \i in {1,...,3}{
    \draw[lien] (\xbus,0 |- S\i.west) -- node[cmt]{Comment ?} (S\i.west);
    \draw[real] (S\i.east) -- (T\i.west);
}
\end{tikzpicture}}%
\caption{FAST de FP2 : les neurones sont répartis sur les c\oe urs, les données restent en VRAM (FC2) et l'orchestrateur lance les calculs à chaque pas.}
\label{fig:fast-fp2}
\end{figure}

\chapter{Schéma synoptique}

\begin{figure}[H]
\centering
\scalebox{0.74}{%
\begin{tikzpicture}[
    >={Stealth[length=2.4mm,width=1.9mm]},
    gbloc/.style={rectangle, rounded corners=5pt, draw=cardline, fill=gray!8, line width=0.9pt,
        align=center, font=\footnotesize\bfseries, text=cardink, inner sep=6pt},
    vrbox/.style={rectangle, rounded corners=5pt, draw=whyborder, fill=whybg, line width=1.1pt,
        align=center, font=\footnotesize\bfseries, text=cardink, inner sep=6pt},
    ent/.style={rectangle, rounded corners=5pt, line width=1pt, align=center,
        font=\footnotesize\bfseries, text=white, inner sep=6pt},
    flux/.style={->, line width=1.1pt, draw=cardline},
    io/.style={<->, line width=1pt, draw=cardline},
    cmd/.style={->, line width=1.1pt, draw=defiborder, dashed},
    flab/.style={font=\scriptsize, text=cardink, fill=white, inner sep=1.5pt, rounded corners=2pt}
]
% Trois colonnes alignées : gauche (x=-5), centre (x=0), droite (x=5)
% --- Rangée du haut -------------------------------------------------------
\node[ent, fill=hubbg, text width=3.6cm, minimum height=1.5cm] (CONF) at (-5,12.0)
    {Configuration\\[-1pt]{\scriptsize modèle, $N$, topologie,\\$dt$, durée, graine}};

\node[gbloc, text width=3.6cm, minimum height=1.5cm] (GEN) at (0,12.0)
    {Générateur de connectome\\[-1pt]{\scriptsize matrice $N \times N$}};

\node[gbloc, text width=3.6cm, minimum height=1.5cm] (REF) at (5,12.0)
    {Référence CPU\\[-1pt]{\scriptsize et temps CPU}};

% --- Milieu : hôte CPU ---------------------------------------------------
\node[gbloc, text width=4.6cm] (ORCH) at (-5,7.8)
    {Orchestrateur de simulation\\[-1pt]{\scriptsize hôte CPU, Python\\-- charge la config et le connectome\\-- lance les calculs\\-- récupère les sorties}};

\node[gbloc, text width=4.6cm] (VALID) at (5,7.3)
    {Validation et mesure\\[-1pt]{\scriptsize -- écart sur $V(t)$\\-- speedup $= t_{CPU} / t_{GPU}$}};

\node[gbloc, text width=4.6cm] (RES) at (5,3.9)
    {Restitution\\[-1pt]{\scriptsize -- fichiers CSV / NumPy, tracés\\-- courbes de temps et speedup}};

% --- Zone GPU --------------------------------------------------------------
\node[vrbox, text width=4.4cm, minimum height=5.4cm] (VRAM) at (-5,-1.6)
    {Mémoire VRAM (résidente)\\[-1pt]{\scriptsize -- états des neurones $V, m, h, n$\\-- connectome dense ou CSR\\-- courants et sorties}};

\node[gbloc, text width=4.4cm, minimum height=1.5cm] (N1) at (4.4,0.1)
    {Courants synaptiques\\[-1pt]{\scriptsize chaque neurone additionne\\ses entrées, selon le connectome}};

\node[gbloc, text width=4.4cm, minimum height=1.5cm] (N2) at (4.4,-3.3)
    {Mise à jour des neurones\\[-1pt]{\scriptsize Hodgkin-Huxley, un fil par neurone}};

\node[font=\footnotesize\bfseries, text=hubbg] (GPUT) at (0,1.8)
    {GPU NVIDIA \textbf{RTX 5090} --- 32~Go de \textbf{VRAM}};

\begin{scope}[on background layer]
\node[rectangle, rounded corners=8pt, draw=hubbg, fill=hubrule!18, line width=1.2pt,
      inner xsep=0.5cm, inner ysep=0.35cm, fit=(VRAM)(N1)(N2)(GPUT)] (GPUB) {};
\end{scope}

% --- Encadrement Epreuvage : reference CPU + validation ---------------------
\begin{scope}[on background layer]
\node[rectangle, rounded corners=8pt, dashed, draw=cardline, line width=1.1pt,
      inner xsep=0.25cm, inner ysep=0.6cm, fit=(REF)(VALID)] (EPR) {};
\end{scope}
\node[anchor=south, font=\scriptsize\bfseries, text=cardink, fill=white, inner sep=2pt]
      at (EPR.north) {Epreuvage};

% --- Flux (tous orthogonaux, aucun croisement) ----------------------------
\draw[flux] (CONF.south) -- node[flab]{configuration} (ORCH.north);
\draw[flux] (GEN.south) |- node[flab, pos=0.3]{matrice $W$} ([yshift=5mm]ORCH.east);
\draw[flux] (REF.south) -- node[flab]{sorties et temps CPU} (VALID.north);
\draw[flux, rounded corners=6pt] (ORCH.east) -- (-1.5,7.8) -- (-1.5,6.35) -- (5,6.35)
    -- node[flab, pos=0.5]{sorties et temps GPU} (RES.north);
\draw[cmd] (ORCH.south) -- node[flab, text=defiborder, pos=0.3]{lancement des calculs} (ORCH.south |- GPUB.north);
\draw[flux] (GPUB.north -| -7.0,0) -- node[flab, pos=0.38]{retour des sorties} (ORCH.south -| -7.0,0);

\draw[io] (VRAM.east |- N1.west) -- node[flab]{lit / écrit} (N1.west);
\draw[io] (VRAM.east |- N2.west) -- node[flab]{lit / écrit} (N2.west);
\draw[flux] (N1.south) -- node[flab]{pas suivant} (N2.north);

% --- Légende --------------------------------------------------------------
\node[draw=cardline, rounded corners=3pt, fill=white, align=left, font=\scriptsize,
      inner sep=4pt, below=0.5cm of GPUB] {%
    \tikz\draw[flux] (0,0)--(0.55,0); \ flux d'information \quad
    \tikz\draw[cmd] (0,0)--(0.55,0); \ commande \quad
    \tikz\fill[whybg, draw=whyborder, line width=0.5pt] (0,0) rectangle (0.28,0.16); \ données gardées sur la carte};
\end{tikzpicture}}%
\caption{Schéma synoptique du POC : l'orchestrateur sur CPU pilote la carte GPU ; la référence CPU et les sorties GPU sont comparées puis restituées.}
\label{fig:syno}
\end{figure}

L'orchestrateur, sur le \textbf{CPU} hôte, charge la configuration et le connectome, lance les calculs à chaque pas et récupère les sorties ; pendant ce temps, toute la simulation vit dans la \textbf{VRAM} de la carte, sans transfert en cours de calcul. Les fonctions cœur se repèrent bloc par bloc : FP1 dans les deux calculs de la carte (courants synaptiques et \textbf{Hodgkin-Huxley}), FP2 dans ces calculs et la mesure du \textbf{speedup}, FC1 dans la validation, FS1 dans le générateur et le stockage dense ou \textbf{CSR}, FC2 dans la carte visée (\textbf{RTX~5090}, 32~Go de \textbf{VRAM}).

